%%%%%%%%%%%%%%%%%%%%%%%%%%%%%%%%%%%%%%%%%%%%%%%%%%%%%%%%%%%%%%%%%%%%%%%%%%%%%%%%%%%%%%%%%%%%%%%%%%%%%%%%%%%%%%%%%%%%%%%%%%%%%%%%%%%%%%%%%%%%%%%%%%%%%%%%%%%
% This is just an example/guide for you to refer to when submitting manuscripts to Frontiers, it is not mandatory to use Frontiers .cls files nor frontiers.tex  %
% This will only generate the Manuscript, the final article will be typeset by Frontiers after acceptance.   
%                                              %
%                                                                                                                                                         %
% When submitting your files, remember to upload this *tex file, the pdf generated with it, the *bib file (if bibliography is not within the *tex) and all the figures.
%%%%%%%%%%%%%%%%%%%%%%%%%%%%%%%%%%%%%%%%%%%%%%%%%%%%%%%%%%%%%%%%%%%%%%%%%%%%%%%%%%%%%%%%%%%%%%%%%%%%%%%%%%%%%%%%%%%%%%%%%%%%%%%%%%%%%%%%%%%%%%%%%%%%%%%%%%%

%%% Version 3.4 Generated 2018/06/15 %%%
%%% You will need to have the following packages installed: datetime, fmtcount, etoolbox, fcprefix, which are normally inlcuded in WinEdt. %%%
%%% In http://www.ctan.org/ you can find the packages and how to install them, if necessary. %%%
%%%  NB logo1.jpg is required in the path in order to correctly compile front page header %%%

\documentclass[utf8]{frontiersSCNS} % for Science, Engineering and Humanities and Social Sciences articles
%\documentclass[utf8]{frontiersHLTH} % for Health articles
%\documentclass[utf8]{frontiersFPHY} % for Physics and Applied Mathematics and Statistics articles

%\setcitestyle{square} % for Physics and Applied Mathematics and Statistics articles
\usepackage{url,hyperref,lineno,microtype,subcaption}
\usepackage[onehalfspacing]{setspace}
\epstopdfsetup{outdir=./}

\linenumbers


% Leave a blank line between paragraphs instead of using \\


\def\keyFont{\fontsize{8}{11}\helveticabold }
\def\firstAuthorLast{Mrad {et~al.}} %use et al only if is more than 1 author
\def\Authors{Assaad Mrad\,$^{1,*}$, Mazen Nakad\,$^{1}$, Yanlan Liu\,$^{1}$ Jean-Christophe Domec\,$^{1,2}$, Sanna Sevanto\,$^{3}$ and Gabriel Katul\,$^{1,4}$}
% Affiliations should be keyed to the author's name with superscript numbers and be listed as follows: Laboratory, Institute, Department, Organization, City, State abbreviation (USA, Canada, Australia), and Country (without detailed address information such as city zip codes or street names).
% If one of the authors has a change of address, list the new address below the correspondence details using a superscript symbol and use the same symbol to indicate the author in the author list.
\def\Address{$^{1}$Nicholas School of the Environment, Duke University, Durham, NC, USA \\
$^{2}$ UMR INRA-ISPA 1391, Bordeaux Sciences Agro, Gradignan 33195, France \\
$^{3}$Earth and Environmental Sciences Division, Los Alamos National Laboratory, Los Alamos, New Mexico, USA \\
$^{4}$Department of Civil and Environmental Engineering, Duke University, Durham, NC}
% The Corresponding Author should be marked with an asterisk
% Provide the exact contact address (this time including street name and city zip code) and email of the corresponding author
\def\corrAuthor{Assaad Mrad}

\def\corrEmail{mradassaad2@gmail.com}




\begin{document}
\onecolumn
\firstpage{1}

\title[Dynamic optimality principle for water use strategies]{A dynamic optimality principle for water use strategies explains isohydric to anisohydric plant responses} 

\author[\firstAuthorLast ]{\Authors} %This field will be automatically populated
\address{} %This field will be automatically populated
\correspondance{} %This field will be automatically populated

\extraAuth{}% If there are more than 1 corresponding author, comment this line and uncomment the next one.
%\extraAuth{corresponding Author2 \\ Laboratory X2, Institute X2, Department X2, Organization X2, Street X2, City X2 , State XX2 (only USA, Canada and Australia), Zip Code2, X2 Country X2, email2@uni2.edu}


\maketitle


\begin{abstract}

%%% Leave the Abstract empty if your article does not require one, please see the Summary Table for full details.
\section{}
For full guidelines regarding your manuscript please refer to \href{http://www.frontiersin.org/about/AuthorGuidelines}{Author Guidelines}.

As a primary goal, the abstract should render the general significance and conceptual advance of the work clearly accessible to a broad readership. References should not be cited in the abstract. Leave the Abstract empty if your article does not require one, please see \href{http://www.frontiersin.org/about/AuthorGuidelines#SummaryTable}{Summary Table} for details according to article type. 


\tiny
 \keyFont{ \section{Keywords:} keyword, keyword, keyword, keyword, keyword, keyword, keyword, keyword} %All article types: you may provide up to 8 keywords; at least 5 are mandatory.
\end{abstract}

\section{Introduction}

Gas exchange at the level of the stomata supports all plant life on earth by providing energy in the form of carbon dioxide ($CO_2$) for growth and maintenance in exchange for water. Over the past millenia, as atmospheric $CO_2$ concentrations decreased ten-fold and an increasing number of plant species adapting to terrestrial ecosystems, the need to optimize carbon gain in relation to water loss arose. This is thought to be behind the adaptation of cacti to arid climates by using a different photosynthetic machinery denoted as Crassulacean Acid Metabolism (CAM) photosynthesis. CAM photosynthesis allows stomata to open at night, when temperature and the vapor pressure deficit (VPD) are low so as to minimize loss of water. It is thought that the stomatal opening of plants utilizing C3 and C4 photosynthesis is also governed by water limitations.

In fact, multiple models grounded on stomatal optimization were put forward and tipped as accurate predictors of transpiration. The accuracy those predictions is crucial for climate forecasting. \textbf{Here offer summary of these models}.

While the assumption that stomatal opening trends are controlled by plant hydraulics is widely used in the plant community, there is a shortage of evidence corroborating this assumption. In this work, the effect of plant water use strategies (from conservative to aggressive) on stomatal conductance ($g$) and leaf water potential ($\psi_l$) is analyzed through a dynamical system perspective that allows setting multiple constraints. 

\citet{Manzoni2013} offer a similar dynamical system approach to the same phenomenon. Water use strategy is prescribed by assigning a 'carbon value' to the terminal soil moisture content around the rooting zone. This study is similarly framed with additional sub-daily measured forcings and constraints on water use such as neighboring tree competition. The results of this model are compared to actual measurements.

The work here answers the following question: to what degree does water use strategy dictate stomatal control during dry-down? In a recent review of the theory of optimal stomatal control, Buckley... recognize the importance of delayed benefits such as keeping soil soil moisture at higher levels for future use. 

\section{Materials and Methods}

\subsection{The model}

The ensure the continuity of the water stream throughout the tree, we model three vertically connected layers: the soil, the trunk and branches, and the leaves. Each of these layers are characterized by a hydraulic resistance that is non-linearly related to water potential. Prescribing hydraulic resistances to these different layers allows computation of quantities regarded as indicative of isohydric to anisohydric behaviour. Such quantities include leaf water potential ($\psi_l$) and stomatal conductance ($g_s$).

At the soil layer, we model the water budget (Proporato...) as follows:
\begin{equation}
    \label{eqn:soil_water}
    \frac{dx(t)}{dt} =\frac{1}{n Z_r}(R - E - L),
\end{equation}
where $x(t)$ is the soil moisture in $m^3m^{-3}$, R, E, and L are the rainfall input into the soil, the evapotranspiration, and uncontrolled losses respectively in $m^2s^{-1}$. $n$ is the soil porosity in $m^3m^{-3}$ and $Z_r$ is considered the plant's rooting depth in $m$. $L$ contains both leakage losses due to soil drainage as well as competition from other plants.

We consider the effect of a drydown on plant stomatal response. During a drydown, there is no rainfall input so $R=0$. The resulting form of Eqn. \ref{eqn:soil_water} is:
\begin{equation}
    \label{eeqn:soil_water_drydown}
    \frac{dx(t)}{dt} = f = - \frac{1}{n Z_r}(E + L),
\end{equation}
where f is the combined losses through evapotranspiration and uncontrolled losses in $s^{-1}$.



To find the optimal trends of $g_s$ and $\psi_l$ with time at different environmental conditions, the method of calculus of variations is used to derive the necessary differential equations (Appendix A). The objective function $J$ in this situation is the following:
\begin{equation}
    \label{eqn:objective}
    J = \int_0^{T_f} A(g_s,T,R_{abs}) dt + J_{T_f},
\end{equation}
where $A$ is the instantaneous plant carbon gain in $molm^{-2}s^{-1}$. $T$ is the atmospheric temperature in $K$, considered equal to the leaf temperature (Appendix B) and $R_{abs}$ is the absorbed radiation in $Wm^{-2}$. $J_{T_f}$ is the carbon value of the terminal soil water moisture $x(T_f)$ prescribed as:
\begin{equation}
    \label{eqn:terminal_gain}
    J_{T_f} = \Lambda x(T_f),
\end{equation}
where $\Lambda$ is a user-defined parameter, in $mol m^{-2}$, prescribing conservative water use strategy when it is large and aggressive water use when it is small. This approach departs from the assumption that carbon gain trends are instantaneously optimal for all plants and. The presence of $J_{T_f}$ in Eqn. \ref{eqn:objective} represents in opportunity cost, measured in carbon gain units, of depleting the soil of water.

The model by Collatz for C3 plants is used to compute $A$ for different plant properties at the leaf level such as Rubisco capacity and for different environmental conditions like temperature and photosynthetically active radiation (PAR). To avoid having to compute derivatives of a piecewise, potentially non-differentiable function, we use the method provided by Campbell and Norman or Giulia Vico.

The effects of high-frequency
fluctuations in light and vapor pressure deficit are neglected,
as they would require modeling of delayed responses of photosynthesis
and conductances (e.g., Vico et al., 2011), a task
beyond the scope here. Limiting the time scale of this investigation
to few weeks also admits another simplification related
to changes in leaf nitrogen status. Increasing leaf nitrogen concentration
improves photosynthetic capacity for a given stomatal
conductance, or allows maintaining relatively high photosynthetic
rates when stomata close under dry conditions (Palmroth et al.,
2013; Weih et al., 2011). Such changes are here assumed to
take place over longer time scales than the typical daily cycles
of light and vapor pressure deficit, but can be accommodated
by varying the kinetic parameters of the photosynthetic demand
curves.

The result is three equations for the soil moisture $x$, $g_s$, and a Lagrange multiplier $\lambda$ commonly referred to as the marginal water use efficiency (MWUE). 

\subsection{Environmental data}

The model required supplying environmental data acting as boundary conditions at the soil and leaf level. These are the initial soil water content ($x_0$), vapor pressure deficit ($VPD$), incident leaf shortwave radiation ($R_l$), and temperature ($T$) (see Table ).

Additionally, different species are characterized by different leaf area index ($LAI$), soil hydraulic conductivity ($g_s$), root-trunk-branch conductivity ($g_t$), rooting depth ($Z_r$). The soil porosity is denoted as $n$.

\section{Results}


\section{Discussion}

\subsection{Equations}
Equations should be inserted in editable format from the equation editor.

\begin{equation}
\sum x+ y =Z\label{eq:01}
\end{equation}

\subsection{Figures}
Frontiers requires figures to be submitted individually, in the same order as they are referred to in the manuscript. Figures will then be automatically embedded at the bottom of the submitted manuscript. Kindly ensure that each table and figure is mentioned in the text and in numerical order. Figures must be of sufficient resolution for publication \href{http://home.frontiersin.org/about/author-guidelines#ResolutionRequirements}{see here for examples and minimum requirements}. Figures which are not according to the guidelines will cause substantial delay during the production process. Please see \href{http://home.frontiersin.org/about/author-guidelines#GeneralStyleGuidelinesforFigures}{here} for full figure guidelines. Cite figures with subfigures as figure \ref{fig:2}B.


\subsubsection{Permission to Reuse and Copyright}
Figures, tables, and images will be published under a Creative Commons CC-BY licence and permission must be obtained for use of copyrighted material from other sources (including re-published/adapted/modified/partial figures and images from the internet). It is the responsibility of the authors to acquire the licenses, to follow any citation instructions requested by third-party rights holders, and cover any supplementary charges.
%%Figures, tables, and images will be published under a Creative Commons CC-BY licence and permission must be obtained for use of copyrighted material from other sources (including re-published/adapted/modified/partial figures and images from the internet). It is the responsibility of the authors to acquire the licenses, to follow any citation instructions requested by third-party rights holders, and cover any supplementary charges.

\subsection{Tables}
Tables should be inserted at the end of the manuscript. Please build your table directly in LaTeX.Tables provided as jpeg/tiff files will not be accepted. Please note that very large tables (covering several pages) cannot be included in the final PDF for reasons of space. These tables will be published as \href{http://home.frontiersin.org/about/author-guidelines#SupplementaryMaterial}{Supplementary Material} on the online article page at the time of acceptance. The author will be notified during the typesetting of the final article if this is the case. 

\section{Nomenclature}

\subsection{Resource Identification Initiative}
To take part in the Resource Identification Initiative, please use the corresponding catalog number and RRID in your current manuscript. For more information about the project and for steps on how to search for an RRID, please click \href{http://www.frontiersin.org/files/pdf/letter_to_author.pdf}{here}.

\subsection{Life Science Identifiers}
Life Science Identifiers (LSIDs) for ZOOBANK registered names or nomenclatural acts should be listed in the manuscript before the keywords. For more information on LSIDs please see \href{http://www.frontiersin.org/about/AuthorGuidelines#InclusionofZoologicalNomenclature}{Inclusion of Zoological Nomenclature} section of the guidelines.


\section{Additional Requirements}

For additional requirements for specific article types and further information please refer to \href{http://www.frontiersin.org/about/AuthorGuidelines#AdditionalRequirements}{Author Guidelines}.

\section*{Conflict of Interest Statement}
%All financial, commercial or other relationships that might be perceived by the academic community as representing a potential conflict of interest must be disclosed. If no such relationship exists, authors will be asked to confirm the following statement: 

The authors declare that the research was conducted in the absence of any commercial or financial relationships that could be construed as a potential conflict of interest.

\section*{Author Contributions}

The Author Contributions section is mandatory for all articles, including articles by sole authors. If an appropriate statement is not provided on submission, a standard one will be inserted during the production process. The Author Contributions statement must describe the contributions of individual authors referred to by their initials and, in doing so, all authors agree to be accountable for the content of the work. Please see  \href{http://home.frontiersin.org/about/author-guidelines#AuthorandContributors}{here} for full authorship criteria.

\section*{Funding}
Details of all funding sources should be provided, including grant numbers if applicable. Please ensure to add all necessary funding information, as after publication this is no longer possible.

\section*{Acknowledgments}
This is a short text to acknowledge the contributions of specific colleagues, institutions, or agencies that aided the efforts of the authors.

\section*{Supplemental Data}
 \href{http://home.frontiersin.org/about/author-guidelines#SupplementaryMaterial}{Supplementary Material} should be uploaded separately on submission, if there are Supplementary Figures, please include the caption in the same file as the figure. LaTeX Supplementary Material templates can be found in the Frontiers LaTeX folder.

\section*{Data Availability Statement}
The datasets [GENERATED/ANALYZED] for this study can be found in the [NAME OF REPOSITORY] [LINK].
% Please see the availability of data guidelines for more information, at https://www.frontiersin.org/about/author-guidelines#AvailabilityofData

\bibliographystyle{frontiersinSCNS_ENG_HUMS} % for Science, Engineering and Humanities and Social Sciences articles, for Humanities and Social Sciences articles please include page numbers in the in-text citations
%\bibliographystyle{frontiersinHLTH&FPHY} % for Health, Physics and Mathematics articles
\bibliography{references}

%%% Make sure to upload the bib file along with the tex file and PDF
%%% Please see the test.bib file for some examples of references

\section*{Figure captions}

%%% Please be aware that for original research articles we only permit a combined number of 15 figures and tables, one figure with multiple subfigures will count as only one figure.
%%% Use this if adding the figures directly in the mansucript, if so, please remember to also upload the files when submitting your article
%%% There is no need for adding the file termination, as long as you indicate where the file is saved. In the examples below the files (logo1.eps and logos.eps) are in the Frontiers LaTeX folder
%%% If using *.tif files convert them to .jpg or .png
%%%  NB logo1.eps is required in the path in order to correctly compile front page header %%%

\begin{figure}[h!]
\begin{center}
\includegraphics[width=10cm]{logo1}% This is a *.eps file
\end{center}
\caption{ Enter the caption for your figure here.  Repeat as  necessary for each of your figures}\label{fig:1}
\end{figure}


\begin{figure}[h!]
\begin{center}
\includegraphics[width=15cm]{logos}
\end{center}
\caption{This is a figure with sub figures, \textbf{(A)} is one logo, \textbf{(B)} is a different logo.}\label{fig:2}
\end{figure}

%%% If you are submitting a figure with subfigures please combine these into one image file with part labels integrated.
%%% If you don't add the figures in the LaTeX files, please upload them when submitting the article.
%%% Frontiers will add the figures at the end of the provisional pdf automatically
%%% The use of LaTeX coding to draw Diagrams/Figures/Structures should be avoided. They should be external callouts including graphics.

\end{document}
